\section{Space Complexity --- Logspace --- Input Access}
%%% generated by the blueprint pipeline from TCSlib.Complexity.SpaceComplexity.Logspace.InputAccess
\subsection*{Space Complexity --- Logspace --- Input Access: Introduction}

Declarations of the Lean module \texttt{TCSlib.Complexity.SpaceComplexity.Logspace.InputAccess}.

\subsection*{Space Complexity --- Logspace --- Input Access: Results}

\begin{theorem}[Nonemptiness of a Pair Encoding]
\lean{Complexity.pairEncode_ne_nil}
\label{Complexity.pairEncode_ne_nil}
\leanok
\uses{Turing.length_pairEncode, Turing.pairEncode}
\difficulty{3}
For any binary strings $x,y$, their pair encoding $\langle x,y\rangle$ is nonempty.
\end{theorem}

\begin{theorem}[Indicator of an encoded index-language word]
\lean{Complexity.indicator_indexLang_pairEncode}
\label{Complexity.indicator_indexLang_pairEncode}
\leanok
\uses{Complexity.LogProg.indicator_eq_decide, Complexity.indexLang, Complexity.pairEncode_mem_indexLang, Turing.MultiTapeTM.indicator, Turing.pairEncode}
\difficulty{3}
For every predicate $p$ on bit strings and natural numbers, bit string $u$, and natural
number $v$, the Boolean indicator of the language of encodings
$	ext{pairEncode}(x,	ext{bits}(i))$ satisfying $p(x,i)$, evaluated at
$	ext{pairEncode}(u,	ext{bits}(v))$, is true exactly when $p(u,v)$ holds and false
otherwise.
\end{theorem}

\begin{theorem}[Whole input with one register]
\lean{Complexity.LogProg.vword_whole₁}
\leanok
\uses{Complexity.LogProg.Mode.seg0, Complexity.LogProg.argSegs, Complexity.LogProg.callSegs, Complexity.LogProg.render, Complexity.LogProg.vword, Turing.pairEncode, Turing.pairEncode_eq_dbl}
\difficulty{1}
For any bit string $u$, decider index $j$, register index $r$, states $l_1,l_0$, and
assignment $V$ of a bit string to each register, the virtual word formed using the
whole-input mode on $u$ and the single register $r$ is the encoded pair $(u,V(r))$.
\end{theorem}

\begin{definition}[Read-phase labels]
\lean{Complexity.LogProg.RdPh}
\label{Complexity.LogProg.RdPh}
\leanok
Define a finite set of read-phase labels consisting of the four parameter-free labels
for equality checking, positive and negative answers, and counter increments, together
with two families indexed by an optional Boolean: two return phases and two clearing
phases.
\end{definition}

\begin{definition}[Read-phase transition function]
\lean{Complexity.LogProg.readTr}
\label{Complexity.LogProg.readTr}
\leanok
\uses{Complexity.LogProg.RdPh, Complexity.LogProg.clrEAct, Complexity.LogProg.eqBAct, Complexity.LogProg.eqCAct, Complexity.LogProg.goAct, Complexity.LogProg.incBAct, Complexity.LogProg.incCAct, Complexity.LogProg.rw2Act, Complexity.LogProg.toEndAct, Complexity.LogProg.xAct, Turing.Action}
For a fixed number $m$ of work tapes, two tape indices $r,s	ext{ in }	ext{Fin}(m)$, a
state assignment from read-phase labels to states, and a state assignment from optional
Boolean values to states, define a transition function. Its input is a read-phase label,
the symbol (or blank) under the selected work-tape head, and the symbols (or blanks)
under all work-tape heads; its output is a one-step action of an $m$-work-tape machine
over the Boolean alphabet. The transition cases implement equality checks and counter
updates, move to the appropriate answer or return phases, and clear the selected
register as specified by the corresponding read phase.
\end{definition}

\begin{definition}[Read-phase machine]
\lean{Complexity.LogProg.readModel}
\label{Complexity.LogProg.readModel}
\leanok
\uses{Complexity.LogProg.RProg, Complexity.LogProg.RdPh, Complexity.LogProg.junkAct, Complexity.LogProg.readTr}
For a natural-number tape count $m$ and two work-tape indices $r,s	ext{ in
}	ext{Fin}(m)$, define an $m$-work-tape machine with Boolean work symbols and states
given by the sum of read-phase labels and optional Boolean values. Its initial state is
the equality-checking phase; in a read-phase state it uses the corresponding read-phase
transition, while in an optional-Boolean state it uses the default action that keeps the
heads stationary and enters no successor state.
\end{definition}

\begin{definition}[Read tape contents at a position]
\lean{Complexity.LogProg.readAt}
\label{Complexity.LogProg.readAt}
\leanok
\uses{Complexity.LogProg.RdPh, Turing.Cfg, Turing.FinTM.bufferTape}
Given a configuration of an $m$-tape Turing machine with Boolean work-tape symbols and
states in the disjoint union of read phases and optional Boolean values, a tape index
$s$, a new state label $l$, a new input-head position $ip$, and a natural number $k$,
define a new configuration by setting its state to $l$, its input-head position to $ip$,
tape $s$ to the tape whose contents encode the binary digits of $k$, and the head
position on tape $s$ to $0$. All other work tapes and head positions, and the output,
are unchanged.
\end{definition}

\begin{definition}[Read-frame configuration relation]
\lean{Complexity.LogProg.readFrame}
\label{Complexity.LogProg.readFrame}
\leanok
\uses{Complexity.LogProg.RdPh, Turing.Cfg}
Let $c$ and $d$ be configurations of an $m$-tape Turing machine with Boolean work-tape
symbols and states in the disjoint union of read phases and optional Boolean values, and
let $r,s$ be work-tape indices. The configuration $d$ is related to $c$ when there are a
state label $l$, an input-head position $ip$, a tape content function $f$, and integers
$p,q$ such that $d$ has state $l$, input-head position $ip$, tape $s$ replaced by $f$,
head positions first updated at $s$ to $p$ and then at $r$ to $q$, and the same output
as $c$; moreover, both $p$ and $q$ lie between $-1$ and the number of binary digits of
the input length plus one, inclusive.
\end{definition}

\begin{theorem}[Register update preserves the read frame]
\lean{Complexity.LogProg.readFrame_reg}
\label{Complexity.LogProg.readFrame_reg}
\leanok
\uses{Complexity.LogProg.RdPh, Complexity.LogProg.readAt, Complexity.LogProg.readFrame, Complexity.LogProg.regCfg, Turing.Cfg}
\difficulty{4}
Let $x$ be a finite Boolean input and let $c$ be a configuration with $m$ work tapes.
Suppose $r$ and $s$ are distinct tape indices and the head on tape $r$ is at position
$0$. For any state labels $l,l'$, valid input-head position $ip$, natural number $k$,
tape content $f:\mathbb Z\to\{\mathrm{none},\mathrm{false},\mathrm{true}\}$, and integer
$p$ with $-1\le p\le \operatorname{length}(	ext{binary representation of }|x|+1)$,
updating tape $s$ and its head to $f$ and $p$ after the read configuration yields a
configuration in the read frame of $c$: the output and other tape contents are
preserved, and the heads on tapes $r,s$ lie between $-1$ and that same length,
inclusive.
\end{theorem}

\begin{theorem}[Read configuration frame]
\lean{Complexity.LogProg.readFrame_at}
\label{Complexity.LogProg.readFrame_at}
\leanok
\uses{Complexity.LogProg.RdPh, Complexity.LogProg.readAt, Complexity.LogProg.readFrame, Complexity.LogProg.readFrame_reg, Complexity.LogProg.regCfg, Turing.Cfg, Turing.FinTM.bufferTape}
\difficulty{4}
Let $x$ be a finite Boolean input, let $c$ be a configuration with $m$ work tapes, and
let $r,s$ be distinct tape indices such that the head on tape $r$ is at position $0$.
For any state label $l$, input-head position $ip$, and natural number $k$, the
configuration obtained from $c$ by entering state $l$, setting the input-head position
to $ip$, and replacing tape $s$ by the binary representation of $k$ with its head at $0$
satisfies the following: its state is present, its input-head position is unrestricted,
its output is unchanged, only tape $s$'s contents are changed, and the head positions
are obtained by changing those on tapes $s$ and $r$; both resulting positions lie
between $-1$ and the length of the binary representation of $|x|+1$, inclusive.
\end{theorem}

\begin{theorem}[Comparison during the read]
\lean{Complexity.LogProg.read_eq}
\label{Complexity.LogProg.read_eq}
\leanok
\uses{Complexity.LogProg.RdPh, Complexity.LogProg.eqCfg, Complexity.LogProg.eq_run, Complexity.LogProg.length_bits_mono, Complexity.LogProg.readAt, Complexity.LogProg.readFrame, Complexity.LogProg.readModel, Complexity.LogProg.rrun, Turing.Cfg, Turing.FinTM.bufferTape}
\difficulty{6}
Let $x$ be a finite Boolean input and let $c$ be a configuration with $m$ work tapes.
Suppose $r$ and $s$ are distinct tape indices, tape $r$ of $c$ contains the binary
representation of $i$, and its head is at position $0$. For any natural number $k\le
|x|$ and valid input-head position $ip$, there is a number $T$ such that the reading
program, started in its initial reading phase with tape $s$ holding the binary
representation of $k$ and input-head position $ip$, reaches the same setup with phase
$yes$ if $k=i$ and phase $no$ otherwise. Every configuration at a time $n<T$ lies in the
read frame of $c$.
\end{theorem}

\begin{theorem}[One step of the read model]
\lean{Complexity.LogProg.read_one}
\label{Complexity.LogProg.read_one}
\leanok
\uses{Complexity.LogProg.RdPh, Complexity.LogProg.inSym, Complexity.LogProg.inputSymbol_eq, Complexity.LogProg.readAt, Complexity.LogProg.readModel, Complexity.LogProg.readTr, Complexity.LogProg.rrun, Complexity.LogProg.rrun_one, Complexity.LogProg.rstep_noncall, Turing.Action.apply, Turing.Cfg, Turing.Cfg.workTapeSymbols, Turing.MultiTapeTM.step}
\difficulty{3}
Let $c$ be a configuration, let $r,s$ be register indices, let $\mathrm{ph}$ be a read
phase, let $\mathrm{ip}$ be an input-head position, and let $k$ be a natural number.
From the configuration obtained by setting the state to $\mathrm{ph}$, the input-head
position to $\mathrm{ip}$, register $s$'s tape to the binary word for $k$, and its head
to $0$, one step of the read model gives the same configuration as applying the action
prescribed by its read transition for input symbol $x[\mathrm{ip}-1]$ (or blank when the
position is $0$ or out of range) and the symbols currently scanned on the work tapes.
\end{theorem}

\begin{theorem}[Cleanup after reading]
\lean{Complexity.LogProg.read_cleanup}
\label{Complexity.LogProg.read_cleanup}
\leanok
\uses{Complexity.LogProg.RdPh, Complexity.LogProg.clr_run, Complexity.LogProg.length_bits_mono, Complexity.LogProg.readAt, Complexity.LogProg.readFrame, Complexity.LogProg.readFrame_at, Complexity.LogProg.readFrame_reg, Complexity.LogProg.readModel, Complexity.LogProg.regCfg, Complexity.LogProg.rewind_x, Complexity.LogProg.rrun, Complexity.LogProg.rrun_add, Complexity.LogProg.xCfg, Turing.Cfg, Turing.FinTM.bufferTape}
\difficulty{6}
Let $x$ be a finite Boolean input and let $c$ be a configuration with $m$ work tapes.
Suppose $r$ and $s$ are distinct tape indices and the head on tape $r$ is at position
$0$. For any natural number $k\le |x|$, valid input-head position $ip$, and optional
Boolean value $a$, there is a number $T$ such that the reading program, started in phase
$rw_1(a)$ with tape $s$ holding the binary representation of $k$ and input-head position
$ip$, reaches the terminal state corresponding to $a$, with input-head position $1$ and
tape $s$'s head at position $0$. Every configuration at a time $n<T$ lies in the read
frame of $c$.
\end{theorem}

\begin{theorem}[One reading advance]
\lean{Complexity.LogProg.read_advance}
\label{Complexity.LogProg.read_advance}
\leanok
\uses{Complexity.LogProg.RdPh, Complexity.LogProg.inSym_succ, Complexity.LogProg.inc_run, Complexity.LogProg.length_bits_mono, Complexity.LogProg.moveInputPos_pos_val, Complexity.LogProg.readAt, Complexity.LogProg.readFrame, Complexity.LogProg.readFrame_at, Complexity.LogProg.readFrame_reg, Complexity.LogProg.readModel, Complexity.LogProg.readTr, Complexity.LogProg.read_eq, Complexity.LogProg.read_one, Complexity.LogProg.regCfg, Complexity.LogProg.rrun, Complexity.LogProg.rrun_add, Complexity.LogProg.rrun_zero, Complexity.LogProg.xAct, Turing.Action.apply, Turing.Cfg, Turing.FinTM.bufferTape}
\difficulty{6}
Let $x$ be a finite Boolean input and let $c$ be a configuration with $m$ work tapes.
Suppose $r$ and $s$ are distinct tape indices, tape $r$ of $c$ contains the binary
representation of $i$, and the head on tape $r$ is at position $0$. If $k<\min(i,|x|)$,
then there is a number $T$ such that the reading program, started in its initial reading
phase with tape $s$ holding the binary representation of $k$ and the input head at
position $k+1$, reaches the same initial reading phase with tape $s$ holding the binary
representation of $k+1$ and the input head at position $k+2$. Every configuration at a
time $n<T$ lies in the read frame of $c$.
\end{theorem}

\begin{theorem}[Completion of the reading routine]
\lean{Complexity.LogProg.read_finish}
\label{Complexity.LogProg.read_finish}
\leanok
\uses{Complexity.LogProg.RdPh, Complexity.LogProg.goAct, Complexity.LogProg.inSym, Complexity.LogProg.inSym_succ, Complexity.LogProg.readAt, Complexity.LogProg.readFrame, Complexity.LogProg.readFrame_at, Complexity.LogProg.readModel, Complexity.LogProg.readTr, Complexity.LogProg.read_cleanup, Complexity.LogProg.read_eq, Complexity.LogProg.read_one, Complexity.LogProg.rrun, Complexity.LogProg.rrun_add, Complexity.LogProg.rrun_zero, Turing.Action.apply, Turing.Cfg, Turing.FinTM.bufferTape}
\difficulty{5}
Let $x$ be a finite Boolean input and let $c$ be a configuration with $m$ work tapes.
Suppose $r$ and $s$ are distinct tape indices, tape $r$ of $c$ contains the binary
representation of $i$, and its head is at position $0$. Then there is a number $T$ such
that the reading program, started in its initial reading phase with tape $s$ holding the
binary representation of $\min(i,|x|)$ and input-head position $\min(i,|x|)+1$, reaches
the terminal state corresponding to the $i$th symbol of $x$ (or no symbol if that index
is out of range), with input-head position $1$ and tape $s$'s head at position $0$.
Every configuration at a time $n<T$ lies in the read frame of $c$.
\end{theorem}

\begin{theorem}[A complete read-model run]
\lean{Complexity.LogProg.read_model_run}
\label{Complexity.LogProg.read_model_run}
\leanok
\uses{Complexity.LogProg.RdPh, Complexity.LogProg.readAt, Complexity.LogProg.readFrame, Complexity.LogProg.readModel, Complexity.LogProg.read_advance, Complexity.LogProg.read_finish, Complexity.LogProg.rrun, Complexity.LogProg.rrun_add, Turing.Cfg, Turing.FinTM.bufferTape}
\difficulty{5}
Let $c$ be a configuration with $m$ registers, let $r\ne s$ be register indices, and let
$i$ be a natural number. Suppose register $r$ contains the binary word for $i$ and its
head is at position $0$. Starting with register $s$ loaded with the binary word for $0$,
its head at $0$, the input head at position $1$, and the state at the initial read
phase, the read model reaches a configuration with register $s$ containing the symbol of
the input at index $i$ (or blank if that index is out of range), its head at $0$, and
the exit state carrying that symbol. There is a natural number $T$ at which this occurs
such that, at every time $n<T$, the configuration has the original output, register
$s$'s tape may have changed, and the heads of registers $s$ and $r$ are at positions $p$
and $q$ satisfying $-1\le p,q\le
\operatorname{length}(\operatorname{bits}(\operatorname{length}(x)+1))$.
\end{theorem}

\begin{theorem}[Read-model and machine runs agree]
\lean{Complexity.LogProg.read_run_eq}
\label{Complexity.LogProg.read_run_eq}
\leanok
\uses{Complexity.LogProg.RdPh, Complexity.LogProg.readModel, Complexity.LogProg.rrun, Complexity.LogProg.rrun_succ, Complexity.LogProg.rstep, Turing.Cfg, Turing.MultiTapeTM.runFrom, Turing.MultiTapeTM.runFrom_succ_eq_step', Turing.MultiTapeTM.step}
\difficulty{4}
For any register indices $r,s$, configuration $c$, and natural number $n$, running the
read model for $n$ steps with the unique empty oracle agrees with running its underlying
multi-tape machine for $n$ steps from $c$.
\end{theorem}

\begin{theorem}[Phases before the read model exits]
\lean{Complexity.LogProg.read_model_phase}
\label{Complexity.LogProg.read_model_phase}
\leanok
\uses{Complexity.LogProg.RdPh, Complexity.LogProg.junkAct, Complexity.LogProg.readModel, Complexity.LogProg.read_run_eq, Complexity.LogProg.rrun, Turing.Action.apply, Turing.Cfg, Turing.MultiTapeTM.runFrom, Turing.MultiTapeTM.runFrom_state_ne_none_of_le, Turing.MultiTapeTM.runFrom_succ_eq_step', Turing.MultiTapeTM.step}
\difficulty{5}
Let $r,s$ be register indices, let $c$ be a configuration, and let $T$ be a natural
number. If the read model is in its exit state carrying
$a\in\{\mathrm{none},\mathrm{false},\mathrm{true}\}$ after $T$ steps from $c$, then at
every time $n<T$ its state is a read phase.
\end{theorem}

\begin{theorem}[State relabeling of reading actions]
\lean{Complexity.LogProg.read_action_map}
\label{Complexity.LogProg.read_action_map}
\leanok
\uses{Complexity.LogProg.RdPh, Complexity.LogProg.clrEAct, Complexity.LogProg.eqBAct, Complexity.LogProg.eqCAct, Complexity.LogProg.goAct, Complexity.LogProg.incBAct, Complexity.LogProg.incCAct, Complexity.LogProg.mv2Act, Complexity.LogProg.readTr, Complexity.LogProg.regAct, Complexity.LogProg.rw2Act, Complexity.LogProg.toEndAct, Complexity.LogProg.xAct, Turing.Action.mapState}
\difficulty{4}
Let $r,s$ be work-tape indices, let $st$ map reading phases to states, and let $ex$ map
optional Boolean values to states. For every reading phase $ph$, optional Boolean value
$a$, and tuple $w$ of symbols under the work-tape heads, mapping the states of the
reading transition defined using the two canonical inclusions into a sum type by the
function that applies $st$ on the left summand and $ex$ on the right gives exactly the
reading transition defined directly using $st$ and $ex$.
\end{theorem}

\begin{theorem}[Execution of the input-reading routine]
\lean{Complexity.LogProg.readIn_run}
\label{Complexity.LogProg.readIn_run}
\leanok
\uses{Complexity.LogProg.RdPh, Complexity.LogProg.readAt, Complexity.LogProg.readModel, Complexity.LogProg.read_action_map, Complexity.LogProg.read_model_phase, Complexity.LogProg.read_model_run, Complexity.LogProg.rrun, Complexity.LogProg.rrun_succ, Complexity.LogProg.rstep_noncall, Turing.Cfg, Turing.Cfg.mapState, Turing.FinTM.bufferTape, Turing.MultiTapeTM.step}
\difficulty{6}
Let $x$ be a finite Boolean input, let $P$ be a register program with $m$ work tapes and
$d$ oracle choices, and let its oracle be a function assigning a Boolean answer to each
oracle choice and Boolean word. Let $r,s$ be distinct tape indices, let $st$ assign a
program state to each reading phase, and let $ex$ assign a program state to each
optional Boolean value. Suppose every transition of the ordinary machine from state
$st(ph)$ is the specified reading transition with register states mapped by $st$ and
terminal states mapped by $ex$, and suppose there are no oracle-call states at any
$st(ph)$. If a configuration $c$ has tape $r$ containing the binary representation of
$i$, tape $s$ empty, both corresponding heads at position $0$, and input-head position
$1$, then some number $T$ of program steps carries the configuration with state $st(eC)$
to the configuration with state $ex$ equal to the $i$th input symbol of $x$ (or no
symbol if that index is out of range). At every time $t<T$, the configuration is in the
read frame of $c$: only tape $s$'s content may change, the output and other tape
contents are preserved, and the heads on tapes $r,s$ lie between $-1$ and the length of
the binary representation of $|x|+1$, inclusive.
\end{theorem}

\begin{definition}[Input-bit language]
\lean{Complexity.inputBitLang}
\label{Complexity.inputBitLang}
\leanok
\uses{Complexity.indexLang}
For a Boolean string $u$, let the predicate at index $i$ hold exactly when the bit of
$u$ at position $i$, using false when $i$ is beyond the end of $u$, is true. The
input-bit language is the indexed language determined by this predicate.
\end{definition}

\begin{definition}[Input-length language]
\lean{Complexity.inputLenLang}
\label{Complexity.inputLenLang}
\leanok
\uses{Complexity.indexLang}
For a Boolean string $u$, let the predicate at index $i$ hold exactly when $i$ is less
than the length of $u$. The input-length language is the indexed language determined by
this predicate.
\end{definition}

\begin{theorem}[Indicator of an input bit]
\lean{Complexity.indicator_inputBitLang}
\label{Complexity.indicator_inputBitLang}
\leanok
\uses{Complexity.indicator_indexLang_pairEncode, Complexity.inputBitLang, Turing.MultiTapeTM.indicator, Turing.pairEncode}
\difficulty{3}
For a binary string $y$ and an index $a\in\mathbb N$, the Boolean indicator of the
language of encoded pairs whose indexed bit is true, evaluated at the encoding of $y$
and the binary representation of $a$, equals the bit of $y$ at index $a$, with
out-of-range indices interpreted as false.
\end{theorem}

\begin{theorem}[Indicator of an input-length query]
\lean{Complexity.indicator_inputLenLang}
\label{Complexity.indicator_inputLenLang}
\leanok
\uses{Complexity.indicator_indexLang_pairEncode, Complexity.inputLenLang, Turing.MultiTapeTM.indicator, Turing.pairEncode}
\difficulty{3}
For a binary string $y$ and an index $a\in\mathbb N$, the Boolean indicator of the
language of encoded pairs whose index is less than the string length, evaluated at the
encoding of $y$ and the binary representation of $a$, equals the Boolean decision of
whether $a<|y|$.
\end{theorem}

\begin{theorem}[Coordinates of a Pair Encoding]
\lean{Complexity.pairEncode_getD}
\label{Complexity.pairEncode_getD}
\leanok
\uses{Turing.dbl, Turing.getElem?_dbl, Turing.length_dbl, Turing.pairEncode, Turing.pairEncode_eq_dbl}
\difficulty{5}
Let $u,v$ be binary strings and let $j\in\mathbb{N}$. The bit at position $j$ of their
pair encoding, taking the default value $0$ when that position is beyond the end, is
$u_{\lfloor j/2\rfloor}$ when $j<2|u|$, is $0$ when $j=2|u|$, is $1$ when $j=2|u|+1$,
and is $v_{j-2|u|-2}$ when $j>2|u|+1$ (with out-of-range bits again taken to be $0$).
\end{theorem}

\begin{theorem}[Second Component of a Pair Encoding]
\lean{Complexity.pairEncode_getD_snd}
\label{Complexity.pairEncode_getD_snd}
\leanok
\uses{Complexity.pairEncode_getD, Turing.pairEncode}
\difficulty{3}
For any binary strings $u,v$ and any $j\in\mathbb{N}$, the bit at position $2|u|+2+j$ of
their pair encoding, taking the default value $0$ when that position is beyond the end,
equals the bit at position $j$ of $v$, with the same default convention.
\end{theorem}

\begin{theorem}[Input length counting loop]
\lean{Complexity.LogProg.inputLength_loop}
\label{Complexity.LogProg.inputLength_loop}
\leanok
\uses{Complexity.LogProg.ARM, Complexity.LogProg.AReach, Complexity.LogProg.AReach.single, Complexity.LogProg.AReach.step, Complexity.LogProg.PreS, Complexity.LogProg.astep, Complexity.LogProg.vword_whole₁, Turing.pairEncode}
\difficulty{6}
Let $A$ be a register program with oracle, and let $r$ be a register and $j$ an oracle
index. Suppose state `scan` calls oracle $j$ on the whole input and the single register
$r$, proceeding to `inc` on a positive answer and `done` otherwise; state `inc`
increments $r$ and returns to `scan`. Suppose that for every $n\le |x|$, the oracle
answers positively on the encoding of $x$ paired with the binary word of $n$ exactly
when $n<|x|$. If $|x|\le B$, all initial registers are at most $B$, and the initial
value of $r$ is at most $|x|$, then the program reaches `done` with $r=|x|$, while every
configuration before that endpoint satisfies $\mathrm{PreS}$ and keeps all registers at
most $B$.
\end{theorem}

\begin{theorem}[Register increment and decrement chains]
\lean{Complexity.LogProg.arm_chain}
\label{Complexity.LogProg.arm_chain}
\leanok
\uses{Complexity.LogProg.ARM, Complexity.LogProg.AReach, Complexity.LogProg.AReach.refl, Complexity.LogProg.AReach.step, Complexity.LogProg.PreS, Complexity.LogProg.astep}
\difficulty{5}
Let $m,k,B\in\mathbb{N}$, let $A$ be a register program with $m$ registers and a state
set $\Lambda$, let $x$ be a bit string, and fix a register $r$, a direction (increment
or decrement), and a sequence of states $(\ell_i)_{i\in\mathbb{N}}$. Suppose that for
each $i<k$, the instruction at state $\ell_i$ increments $r$ and goes to $\ell_{i+1}$ in
the increment case, or decrements $r$ and goes to $\ell_{i+1}$ in the decrement case. If
the initial register values $v$ satisfy $v(s)\le B$ for every register $s$, and
additionally $v(r)+k\le B$ in the increment case or $k\le v(r)$ in the decrement case,
then the program reaches the configuration with state $\ell_k$, unchanged result field,
and register values given by replacing $v(r)$ by $v(r)+k$ or by the truncated difference
$v(r)-k$, respectively. Throughout this finite run, every configuration satisfies the
program's side conditions and every register value is at most $B$.
\end{theorem}

\begin{lemma}[Repeated increments preserve bounded reachability]
\lean{Complexity.LogProg.arm_increments}
\label{Complexity.LogProg.arm_increments}
\leanok
\uses{Complexity.LogProg.ARM, Complexity.LogProg.AReach, Complexity.LogProg.PreS, Complexity.LogProg.arm_chain}
\difficulty{1}
Let $A$ be an arithmetic register program, let $x$ be a Boolean input, and let $r$ be
one of its registers. Suppose $k,B\in\mathbb{N}$, the instruction at each of the states
$\ell_0,\ldots,\ell_{k-1}$ increments register $r$ and proceeds to the next state, and
an initial register valuation $v$ satisfies $v(s)\le B$ for every register $s$ and
$v(r)+k\le B$. Then the program can reach the configuration with state $\ell_k$,
register valuation obtained by replacing $v(r)$ with $v(r)+k$, and no answer yet, from
the configuration with state $\ell_0$, valuation $v$, and no answer yet, with the
precondition holding at every preceding step and all register values bounded by $B$.
\end{lemma}

\begin{theorem}[Aligned-pair scan loop]
\lean{Complexity.LogProg.alignedPair_loop}
\label{Complexity.LogProg.alignedPair_loop}
\leanok
\uses{Complexity.LogProg.AConf, Complexity.LogProg.ARM, Complexity.LogProg.AReach, Complexity.LogProg.AReach.single, Complexity.LogProg.AReach.step, Complexity.LogProg.AReach.trans, Complexity.LogProg.PreS, Complexity.LogProg.arm_increments, Complexity.LogProg.astep, Complexity.LogProg.vword_whole₁, Turing.pairDecode, Turing.pairEncode}
\difficulty{7}
Let $A$ be a register program with oracle, let $x,w$ be finite bit strings, and let $r$
and $\ell$ be distinct registers. Choose an oracle index $j$, program states
$\mathrm{scan},\mathrm{first},\mathrm{fail},\mathrm{done}$ and states
$\mathrm{test}(b),\mathrm{read}(b)$ for each bit $b$, and states $\mathrm{advance}(k,i)$
for $k\in\{0,1,2,3\}$ and natural numbers $i$. Suppose the scan instruction compares
registers $r$ and $\ell$, the first instruction calls oracle $j$ on the whole input and
the word in register $r$, each test compares those registers, and each read calls the
oracle on that same input, branching as specified by the bit being tested; each advance
instruction increments register $r$ and proceeds to the next advance state. Assume the
four advance states at offset $\mathrm{stride}$ lead respectively to
$\mathrm{test}(\mathrm{false})$, $\mathrm{test}(\mathrm{true})$, $\mathrm{scan}$, and
$\mathrm{done}$; also assume $0<\mathrm{stride}$, every register value in the initial
vector $v$ is at most $B$, $v(\ell)=v(r)+\mathrm{stride}|w|$, and for every $i<|w|$ the
oracle answer on $\langle x,\mathrm{bits}(v(r)+\mathrm{stride}i)\rangle$ is the $i$-th
bit of $w$ (defaulting to false past the end). Then the program reaches a failure state
with register $r$ increased by $\mathrm{stride}t$ for some $t\le |w|$ if the
aligned-pair parser rejects $w$; if it parses $w$ as $(a,b)$, it reaches $\mathrm{done}$
with register $r$ increased by $\mathrm{stride}(2|a|+2)$. In either case, throughout the
finite run every intermediate configuration satisfies the program's
precondition—equality tests use distinct registers, input-word tests occur on valid
encodings, and call argument lists have no repetitions—and every register remains at
most $B$.
\end{theorem}

\begin{theorem}[Register update preserves an upper bound]
\lean{Complexity.LogProg.register_update_le}
\label{Complexity.LogProg.register_update_le}
\leanok
\difficulty{4}
Let $v$ assign a natural-number value to each of $m$ registers, let $r$ be a register,
and let $B,n$ be natural numbers. If every value assigned by $v$ is at most $B$ and
$n\leq B$, then replacing the value at $r$ by $n$ leaves every register value at most
$B$.
\end{theorem}

\begin{theorem}[Copying one register by a loop]
\lean{Complexity.LogProg.registerCopy_loop}
\label{Complexity.LogProg.registerCopy_loop}
\leanok
\uses{Complexity.LogProg.ARM, Complexity.LogProg.AReach, Complexity.LogProg.AReach.single, Complexity.LogProg.AReach.step, Complexity.LogProg.PreS, Complexity.LogProg.astep, Complexity.LogProg.register_update_le}
\difficulty{5}
Let a register program $A$ have distinct destination and source registers, and let its
instruction at state $test$ compare these registers, proceeding to $done$ when they are
equal and to $bump$ otherwise; at $bump$ it increments the destination and returns to
$test$. Suppose every register initially contains a value at most $B$, and the
destination's initial value is at most the source's. From the configuration at $test$,
with no result yet, the program can reach $done$ with the destination set to the
source's initial value and all other registers unchanged, while at every configuration
before the final one the precondition invariant holds and all register values remain at
most $B$.
\end{theorem}

\begin{theorem}[Scaling one register by a loop]
\lean{Complexity.LogProg.registerScale_loop}
\label{Complexity.LogProg.registerScale_loop}
\leanok
\uses{Complexity.LogProg.ARM, Complexity.LogProg.AReach, Complexity.LogProg.AReach.single, Complexity.LogProg.AReach.step, Complexity.LogProg.AReach.trans, Complexity.LogProg.PreS, Complexity.LogProg.arm_increments, Complexity.LogProg.astep, Complexity.LogProg.register_update_le}
\difficulty{6}
Let a register program $A$ have distinct source and destination registers. At state
$test$ it branches to $done$ if the source is zero and otherwise to $down$; at $down$ it
decrements the source and enters $inc(0)$. For every $i<factor$, state $inc(i)$
increments the destination and enters $inc(i+1)$, and $inc(factor)=test$. If every
initial register value is at most $B$ and the initial destination value plus $factor$
times the initial source value is at most $B$, then from the configuration at $test$ the
program can reach $done$ with the source set to zero and the destination increased by
$factor$ times the initial source value. Throughout the configurations before the final
one, the precondition invariant holds and every register value is at most $B$.
\end{theorem}

\begin{lemma}[Binary width after halving]
\lean{Complexity.LogProg.access_width_half}
\label{Complexity.LogProg.access_width_half}
\leanok
\uses{Complexity.LogProg.bits_half}
\difficulty{3}
For every nonzero natural number $n$, the binary word for $\lfloor n/2\rfloor$ has
length one less than the binary word for $n$.
\end{lemma}

\begin{theorem}[Canonical words and final bits]
\lean{Complexity.LogProg.drop_canon}
\label{Complexity.LogProg.drop_canon}
\leanok
\uses{Complexity.LogProg.Canon}
\difficulty{4}
For every Boolean word $w$, its last bit with the default value true is true if and only
if either $w$ is empty or its last bit is true.
\end{theorem}

\begin{theorem}[Canonical tail loop]
\lean{Complexity.LogProg.canonicalTail_loop}
\label{Complexity.LogProg.canonicalTail_loop}
\leanok
\uses{Complexity.LogProg.ARM, Complexity.LogProg.AReach, Complexity.LogProg.AReach.single, Complexity.LogProg.AReach.step, Complexity.LogProg.AReach.trans, Complexity.LogProg.PreS, Complexity.LogProg.access_width_half, Complexity.LogProg.arm_increments, Complexity.LogProg.astep, Complexity.LogProg.register_update_le, Complexity.LogProg.vword_whole₁, Turing.pairEncode}
\difficulty{6}
Let $A$ be a program with $m$ registers, control states in $\Lambda$, and oracle
deciders indexed by $d$, and choose pairwise distinct registers
$\mathsf{ptr},\mathsf{limit},\mathsf{cap}$. Suppose the loop compares the pointer with
the limit, reads a decider answer at the binary encoding of the pointer value, halves
the cap register, and advances the pointer by a fixed stride before repeating; on
reaching the limit it reports done or fail according to a supplied final bit, while a
zero cap or a failed length check leads to fail. Assume $\mathsf{stride}>0$ and the
stated transitions of $A$ implement these branches, query, halving, and stride
increments exactly. Let $w$ be a Boolean word, let $\mathit{last}$ be a Boolean bit, and
suppose all initial registers are at most $B$, the limit equals the pointer plus
$\mathsf{stride}\,|w|$, and each oracle query at pointer value plus $\mathsf{stride}\,i$
returns bit $i$ of $w$ for every $i<|w|$ (defaulting to false when absent). Then there
are final register values, still bounded by $B$, with the pointer at most the limit and
every register other than the pointer and cap unchanged, such that the program reaches
the control state determined as follows: if $|w|$ is at most the bit length of the
initial cap, it reports done when the last bit of $w$ (defaulting to $\mathit{last}$ for
an empty word) is true, and fail otherwise; if the length exceeds that bit length, it
reports fail. Every intervening configuration satisfies the program’s safety conditions
and keeps all registers bounded by $B$.
\end{theorem}

\begin{theorem}[Bounded decoding stride loop]
\lean{Complexity.LogProg.boundedDecodeStride_loop}
\label{Complexity.LogProg.boundedDecodeStride_loop}
\leanok
\uses{Complexity.LogProg.ARM, Complexity.LogProg.AReach, Complexity.LogProg.AReach.single, Complexity.LogProg.AReach.step, Complexity.LogProg.AReach.trans, Complexity.LogProg.PreS, Complexity.LogProg.arm_chain, Complexity.LogProg.astep, Complexity.LogProg.bitsVal, Complexity.LogProg.registerCopy_loop, Complexity.LogProg.registerScale_loop, Complexity.LogProg.register_update_le, Complexity.LogProg.vword_whole₁, Turing.pairEncode}
\difficulty{7}
Let $m,d$ be natural numbers, let $\Lambda$ be a type of machine states, let $A$ assign
an instruction to each state, and let an oracle assign a Boolean answer to each index in
$\{0,\ldots,d-1\}$ and each bit string. Fix an input bit string $x$, four pairwise
distinct register indices $\mathrm{ptr},\mathrm{base},\mathrm{acc},\mathrm{tmp}$ among
the $m$ registers, an oracle index $j<d$, states
$\mathrm{test},\mathrm{read},\mathrm{done}$ and state families $\mathrm{downs}(i)$,
$\mathrm{copy}(b)$, $\mathrm{bump}(b)$, $\mathrm{clear}(b)$, $\mathrm{scale}(b)$,
$\mathrm{take}(b)$, $\mathrm{add}(b)$, and $\mathrm{incs}(b,i)$, and a positive stride
$s$. Suppose the instruction at $\mathrm{test}$ compares the pointer and base registers,
proceeding to $\mathrm{done}$ if equal and otherwise to $\mathrm{downs}(0)$; for each
$i<s$, $\mathrm{downs}(i)$ decrements the pointer, and
$\mathrm{downs}(s)=\mathrm{read}$. The state $\mathrm{read}$ calls oracle $j$ in
whole-input mode with the pointer register as its sole argument, proceeding to
$\mathrm{copy}(1)$ on a positive answer and $\mathrm{copy}(0)$ otherwise. For either bit
$b$, $\mathrm{copy}(b)$ compares the temporary and accumulator registers, proceeding to
$\mathrm{clear}(b)$ if equal and $\mathrm{bump}(b)$ otherwise; $\mathrm{bump}(b)$
increments the temporary register and returns to $\mathrm{copy}(b)$; $\mathrm{clear}(b)$
clears the accumulator and proceeds to $\mathrm{scale}(b)$; $\mathrm{scale}(b)$ proceeds
to $\mathrm{add}(b)$ when the temporary register is zero and to $\mathrm{take}(b)$
otherwise; $\mathrm{take}(b)$ decrements the temporary register and proceeds to
$\mathrm{incs}(b,0)$. For $i<2$, $\mathrm{incs}(b,i)$ increments the accumulator and
proceeds to $\mathrm{incs}(b,i+1)$, with $\mathrm{incs}(b,2)=\mathrm{scale}(b)$. The
state $\mathrm{add}(b)$ increments the accumulator and returns to $\mathrm{test}$ if
$b=1$, and otherwise proceeds to $\mathrm{test}$ if the temporary register is zero (or
remains at $\mathrm{add}(b)$ if it is not). Let $w$ be a bit string, $B,k$ natural
numbers, and let $v$ assign a natural number to each register. Assume that for every
$i\le |w|$, the value of the binary suffix of $w$ starting at position $i$ is at most
$B$; $k\le |w|$; every register value in $v$ is at most $B$; the pointer value is
$v(\mathrm{base})+sk$, the accumulator value is the value of that suffix starting at
$k$, and the temporary register is zero. Also suppose that for every $i<|w|$, the oracle
answer on the encoding of $x$ paired with the binary expansion of $v(\mathrm{base})+si$
equals the $i$th bit of $w$ (using $0$ when that position is beyond the list). Then
there is a register assignment $v'$ with all register values at most $B$, pointer value
$v'(\mathrm{ptr})=v(\mathrm{base})$, accumulator value equal to the value of the entire
binary list $w$, and temporary register zero; all other registers retain their values.
Moreover, the machine reaches configuration $(\mathrm{done},v',\mathrm{none})$ from
$(\mathrm{test},v,\mathrm{none})$ in finitely many steps, and at every step before
arrival the configuration satisfies the stated precondition and all register values
remain at most $B$.
\end{theorem}

\begin{theorem}[Bounded decoding loop]
\lean{Complexity.LogProg.boundedDecode_loop}
\label{Complexity.LogProg.boundedDecode_loop}
\leanok
\uses{Complexity.LogProg.ARM, Complexity.LogProg.AReach, Complexity.LogProg.PreS, Complexity.LogProg.bitsVal, Complexity.LogProg.boundedDecodeStride_loop, Turing.pairEncode}
\difficulty{3}
Let $A$ be a finite-register program with $m$ registers and control states in $\Lambda$,
and let an oracle assign a Boolean answer to each of $d$ deciders on every bit string.
Choose pairwise distinct registers
$\mathsf{ptr},\mathsf{base},\mathsf{acc},\mathsf{tmp}$ and control states implementing
the following loop: compare $\mathsf{ptr}$ with $\mathsf{base}$, decrement the pointer
if they differ, query decider $j$ on the doubled input paired with the binary
representation of the pointer value, and use the answer to either clear the accumulator
or increment the temporary register. Then, if the temporary register is zero, add the
answer bit to the accumulator (or leave it unchanged for a negative answer); otherwise
decrement the temporary register and increment the accumulator twice. Assume the
indicated transitions of $A$ implement precisely these instructions, with the two-step
increment chain ending at the scaling state. For a word $w$, let $\operatorname{val}(w)$
be its binary value with the first bit least significant. Suppose $B$ bounds the value
of every suffix of $w$, the initial register values are all at most $B$, the pointer is
$k$ above the base, the accumulator equals $\operatorname{val}(w[k:])$, the temporary
register is zero, and every oracle query at base value plus $i$ returns the $i$-th bit
of $w$ (taking a missing bit to be zero). Then there are final register values, still
bounded by $B$, at which the pointer equals the base, the accumulator equals
$\operatorname{val}(w)$, the temporary register is zero, and every other register is
unchanged. Moreover, the execution reaches the stated done control state from the
specified initial configuration, with every intervening configuration satisfying the
program’s safety conditions (distinct registers for equality tests, valid input for
input comparisons, distinct arguments for oracle calls) and keeping all registers
bounded by $B$.
\end{theorem}

\begin{theorem}[Indexed seek loop]
\lean{Complexity.LogProg.indexedSeek_loop}
\label{Complexity.LogProg.indexedSeek_loop}
\leanok
\uses{Complexity.LogProg.ARM, Complexity.LogProg.AReach, Complexity.LogProg.AReach.single, Complexity.LogProg.AReach.step, Complexity.LogProg.AReach.trans, Complexity.LogProg.PreS, Complexity.LogProg.arm_increments, Complexity.LogProg.astep, Complexity.LogProg.register_update_le}
\difficulty{6}
Let $A$ be a register program with oracle, and let
$\mathrm{ptr},\mathrm{limit},\mathrm{idx},\mathrm{query}$ be pairwise distinct
registers. Suppose the program states implement the following loop: at `test`, compare
$\mathrm{ptr}$ with $\mathrm{limit}$ and proceed to `miss` if equal, otherwise to
`probe`; at `probe`, compare $\mathrm{idx}$ with $\mathrm{query}$ and proceed to `hit`
if equal, otherwise to `bump`; at `bump`, increment $\mathrm{idx}$ and enter
$\mathrm{advance}(0)$; each state $\mathrm{advance}(i)$ for $i<\mathrm{stride}$
increments $\mathrm{ptr}$ and advances to $\mathrm{advance}(i+1)$, with
$\mathrm{advance}(\mathrm{stride})=\mathrm{test}$. Assume $0<\mathrm{stride}$, every
register initially contains at most $B$,
$\mathrm{limit}=\mathrm{ptr}+\mathrm{stride}\cdot\mathrm{count}$,
$\mathrm{idx}+\mathrm{count}\le B$, and $\mathrm{idx}\le\mathrm{query}$. If
$\mathrm{query}<\mathrm{idx}+\mathrm{count}$, the run reaches `hit` with $\mathrm{ptr}$
increased by $\mathrm{stride}(\mathrm{query}-\mathrm{idx})$ and $\mathrm{idx}$ set to
$\mathrm{query}$; otherwise it reaches `miss` with $\mathrm{ptr}$ set to its initial
$\mathrm{limit}$ and $\mathrm{idx}$ increased by $\mathrm{count}$. In either case, every
configuration before the endpoint satisfies the program's precondition $\mathrm{PreS}$
and keeps every register at most $B$.
\end{theorem}

\begin{definition}[Access labels]
\lean{Complexity.LogProg.AccessLabel}
\label{Complexity.LogProg.AccessLabel}
\leanok
Define access labels to be the finite collection consisting of a scan label; a second
label carrying a Boolean value; a tail label carrying a Boolean value and an optional
Boolean value; a data label carrying a Boolean value; and a second-data label carrying
two Boolean values.
\end{definition}

\begin{definition}[Access states]
\lean{Complexity.LogProg.AccessState}
\label{Complexity.LogProg.AccessState}
\leanok
\uses{Complexity.LogProg.AccessLabel, Complexity.LogProg.RdPh}
Define access states to include read states carrying an access label and a read phase;
received-value states carrying an access label and an optional Boolean value; increment
states carrying one of four register indices, an access label, and a Boolean value;
restore states carrying an optional Boolean value; clear states carrying an optional
Boolean value and a Boolean value; bump states carrying a Boolean value; and answer
states carrying a Boolean value.
\end{definition}

\begin{definition}[Register selected by an access label]
\lean{Complexity.LogProg.accessReg}
\label{Complexity.LogProg.accessReg}
\leanok
\uses{Complexity.LogProg.AccessLabel}
Assign register $3$ to each data label and each second-data label, and assign register
$0$ to every other access label.
\end{definition}

\begin{definition}[Access state after a read]
\lean{Complexity.LogProg.accessExit}
\label{Complexity.LogProg.accessExit}
\leanok
\uses{Complexity.LogProg.AccessState}
For an optional Boolean value $a$, define the resulting access state to be the answer
state carrying false when $a$ is absent, and otherwise the read state carrying the data
label for $a$ and the initial read phase.
\end{definition}

\begin{definition}[Access transition function]
\lean{Complexity.LogProg.accessTr}
\label{Complexity.LogProg.accessTr}
\leanok
\uses{Complexity.LogProg.AccessState, Complexity.LogProg.accessExit, Complexity.LogProg.accessReg, Complexity.LogProg.backAct, Complexity.LogProg.clrEAct, Complexity.LogProg.goAct, Complexity.LogProg.incBAct, Complexity.LogProg.incCAct, Complexity.LogProg.readTr, Complexity.LogProg.retAct, Complexity.LogProg.toEndAct, Complexity.LogProg.xAct, Turing.Action}
For a Boolean mode parameter, define the transition function of the four-work-tape
access procedure on its access states by the specified cases: it performs the read,
comparison, counter-update, restoration, and answer steps, and returns the corresponding
one-step machine action.
\end{definition}

\begin{definition}[Access program]
\lean{Complexity.LogProg.accessProg}
\label{Complexity.LogProg.accessProg}
\leanok
\uses{Complexity.LogProg.AccessState, Complexity.LogProg.RProg, Complexity.LogProg.accessTr}
For each Boolean mode, define a deterministic Turing program with four work tapes and no
deciders. Its initial state is the state that reads the scan label in phase $eC$, its
transition function is the access transition function for that mode, and every state has
no associated call specification.
\end{definition}

\begin{definition}[Access configuration]
\lean{Complexity.LogProg.accessCfg}
\label{Complexity.LogProg.accessCfg}
\leanok
\uses{Complexity.LogProg.AccessState, Turing.Cfg, Turing.FinTM.bufferTape}
Given a Boolean input $y$, an access state $s$, and natural numbers $i,k,p$, define a
four-work-tape configuration on input $y$ with state $s$, input-head position $1$, empty
output, and work-tape contents indexed by $r\in\{0,1,2,3\}$ equal to the buffered binary
representation of the $r$th entry of $(i,0,k,p)$. The head on work tape $2$ is at
position $q$ and the other work-tape heads are at position $0$, where $q$ defaults to
$0$.
\end{definition}

\begin{definition}[Safe access configuration]
\lean{Complexity.LogProg.accessSafe}
\label{Complexity.LogProg.accessSafe}
\leanok
\uses{Complexity.LogProg.AccessState, Turing.Cfg}
For a Boolean input $y$, a four-work-tape configuration $c$ on input $y$ is safe if
every work-tape head position lies between $-1$ and the length of the binary
representation of $|y|+1$, inclusive.
\end{definition}

\begin{definition}[Safe access path]
\lean{Complexity.LogProg.AccessPath}
\label{Complexity.LogProg.AccessPath}
\leanok
\uses{Complexity.LogProg.AccessState, Complexity.LogProg.Dbl.o, Complexity.LogProg.accessProg, Complexity.LogProg.accessSafe, Complexity.LogProg.rstep, Turing.Cfg}
For a Boolean mode and a Boolean input $y$, define an access path to be a reflexive
transitive sequence of four-work-tape configurations on input $y$ in which each step
goes from a safe configuration to the result of one transition of the access program for
that mode, using the doubled-input convention.
\end{definition}

\begin{lemma}[Access paths are safe run segments]
\lean{Complexity.LogProg.accessPath_run}
\label{Complexity.LogProg.accessPath_run}
\leanok
\uses{Complexity.LogProg.AccessPath, Complexity.LogProg.AccessState, Complexity.LogProg.Dbl.o, Complexity.LogProg.accessProg, Complexity.LogProg.accessSafe, Complexity.LogProg.rrun, Complexity.LogProg.rrun_succ, Turing.Cfg}
\difficulty{4}
For either access mode, if there is an access path from a configuration $c$ on input $y$
to a configuration $c'$, then some natural number $T$ satisfies that the run of the
access program from $c$ reaches $c'$ at time $T$, and every configuration at times $t<T$
has all work-head positions between $-1$ and the length of the binary representation of
$|y|+1$, inclusive.
\end{lemma}

\begin{lemma}[Safe runs define access paths]
\lean{Complexity.LogProg.accessPath_of_run}
\label{Complexity.LogProg.accessPath_of_run}
\leanok
\uses{Complexity.LogProg.AccessPath, Complexity.LogProg.AccessState, Complexity.LogProg.Dbl.o, Complexity.LogProg.accessProg, Complexity.LogProg.accessSafe, Complexity.LogProg.rrun, Complexity.LogProg.rrun_succ, Turing.Cfg}
\difficulty{3}
For either access mode, let $c$ be a configuration on input $y$, and let $T$ be a
natural number. If every configuration at times $t<T$ in the run of the access program
from $c$ has each work head between $-1$ and the length of the binary representation of
$|y|+1$, inclusive, then the run from $c$ to time $T$ is an access path. Here an access
path is a finite sequence of zero or more program steps, with each step allowed when its
starting configuration satisfies that same head-position bound.
\end{lemma}

\begin{lemma}[Safety of an Access Configuration]
\lean{Complexity.LogProg.accessCfg_safe}
\label{Complexity.LogProg.accessCfg_safe}
\leanok
\uses{Complexity.LogProg.AccessState, Complexity.LogProg.accessCfg, Complexity.LogProg.accessSafe}
\difficulty{3}
Let $y$ be a Boolean word, let $st$ be an access state, let $i,k,p$ be natural numbers,
and let $q$ be an integer satisfying $-1\le q\le |\mathrm{bits}(|y|+1)|$. Consider the
access configuration whose state is $st$, whose input-head position is $1$, whose four
work tapes store the binary expansions of the corresponding entries of $(i,0,k,p)$,
whose work-head positions are $(0,0,q,0)$, and whose output is empty. Each of its
work-head positions lies between $-1$ and the length of the binary expansion of $|y|+1$,
inclusive.
\end{lemma}

\begin{lemma}[Reading an input bit]
\lean{Complexity.LogProg.access_read}
\label{Complexity.LogProg.access_read}
\leanok
\uses{Complexity.LogProg.AccessLabel, Complexity.LogProg.AccessPath, Complexity.LogProg.Dbl.o, Complexity.LogProg.accessCfg, Complexity.LogProg.accessCfg_safe, Complexity.LogProg.accessPath_of_run, Complexity.LogProg.accessProg, Complexity.LogProg.accessReg, Complexity.LogProg.readIn_run, Complexity.LogProg.rrun}
\difficulty{5}
Let $y$ be a Boolean word, let $l$ be an access label, and let $i,k,p$ be natural
numbers. If the work-tape head position $q$ satisfies $-1\le q\le
|\mathrm{bits}(|y|+1)|$, then there is a finite safe computation path from the
configuration in read phase $eC$ with label $l$, parameters $i,k,p$, and head position
$q$ to the configuration in which the read result is the optional bit of $y$ at index
$i$ for ordinary labels or index $p$ for data labels, with the same other parameters and
head position.
\end{lemma}

\begin{lemma}[Preservation of Head Bounds]
\lean{Complexity.LogProg.access_frame}
\label{Complexity.LogProg.access_frame}
\leanok
\uses{Complexity.LogProg.AccessState, Complexity.LogProg.accessSafe, Complexity.LogProg.regCfg, Turing.Cfg}
\difficulty{4}
Let $y$ be a Boolean input word and $c$ a configuration of the four-tape access machine
on input $y$, with every work-tape head between $-1$ and the bit-length of $|y|+1$. If
$r$ is one of the four work tapes, $s$ is any access-machine state, $f$ is any tape
content, and $q$ lies between $-1$ and the bit-length of $|y|+1$, then the configuration
obtained from $c$ by replacing tape $r$ with $f$, placing its head at $q$, and setting
the state to $s$ still has every work-tape head within those bounds.
\end{lemma}

\begin{lemma}[Incrementing a Binary Counter]
\lean{Complexity.LogProg.access_inc}
\label{Complexity.LogProg.access_inc}
\leanok
\uses{Complexity.LogProg.AccessLabel, Complexity.LogProg.AccessPath, Complexity.LogProg.AccessState, Complexity.LogProg.Dbl.o, Complexity.LogProg.accessPath_of_run, Complexity.LogProg.accessProg, Complexity.LogProg.accessSafe, Complexity.LogProg.access_frame, Complexity.LogProg.inc_run, Complexity.LogProg.length_bits_mono, Complexity.LogProg.regCfg, Turing.Cfg, Turing.FinTM.bufferTape}
\difficulty{4}
Let $m$ be a Boolean mode, $y$ a Boolean input word, $r$ one of the four work tapes,
$\ell$ an access label, and $c$ a configuration of the four-tape access machine on input
$y$. Suppose every work-tape head of $c$ is between $-1$ and the bit-length of $|y|+1$,
and $n+1\le |y|+1$. There is a finite sequence of access-machine steps from the
configuration obtained by replacing tape $r$ with the binary representation of $n$,
placing its head at $0$, and setting the state to “increment on tape $r$ with label
$\ell$,” to the configuration obtained by replacing that tape with the binary
representation of $n+1$, placing its head at $0$, and setting the state to “read label
$\ell$ in phase eC”; all other configuration data are retained.
\end{lemma}

\begin{lemma}[Clearing a Work Tape]
\lean{Complexity.LogProg.access_clear}
\label{Complexity.LogProg.access_clear}
\leanok
\uses{Complexity.LogProg.AccessPath, Complexity.LogProg.AccessState, Complexity.LogProg.Dbl.o, Complexity.LogProg.accessExit, Complexity.LogProg.accessPath_of_run, Complexity.LogProg.accessProg, Complexity.LogProg.accessSafe, Complexity.LogProg.access_frame, Complexity.LogProg.clr_run, Complexity.LogProg.length_bits_mono, Complexity.LogProg.regCfg, Turing.Cfg, Turing.FinTM.bufferTape}
\difficulty{4}
Let $y$ be a Boolean input word, let $m$ be a Boolean mode, let $e$ be an optional
Boolean, and let $c$ be a configuration of the four-tape access machine on input $y$.
Suppose every work-tape head of $c$ lies between $-1$ and the bit-length of $|y|+1$, and
let $n\le |y|+1$. There is a finite sequence of access-machine steps, each taken from a
configuration satisfying these head-position bounds, from the configuration obtained by
setting tape $0$ to the binary word for $n$, placing its head at $0$, and setting the
state to the clearing state with parameter $e$, to the configuration obtained by setting
tape $0$ to blank, keeping its head at $0$, and setting the state to the exit state
determined by $e$; all other configuration data are retained.
\end{lemma}

\begin{lemma}[One access-machine transition]
\lean{Complexity.LogProg.access_step}
\label{Complexity.LogProg.access_step}
\leanok
\uses{Complexity.LogProg.AccessPath, Complexity.LogProg.AccessState, Complexity.LogProg.Dbl.o, Complexity.LogProg.accessProg, Complexity.LogProg.accessSafe, Complexity.LogProg.accessTr, Complexity.LogProg.rstep_noncall, Turing.Action.apply, Turing.Cfg, Turing.Cfg.inputSymbol, Turing.Cfg.workTapeSymbols, Turing.MultiTapeTM.step}
\difficulty{3}
Let $y$ be a Boolean word, $c$ a configuration, and $st$ an access state. If the state
of $c$ is $st$ and every work-tape head position of $c$ lies between $-1$ and
$|\mathrm{bits}(|y|+1)|$, inclusive, then there is a finite safe computation path from
$c$ to the configuration obtained by applying the access-machine transition for mode and
state $st$ to $c$.
\end{lemma}

\begin{lemma}[Incrementing a Counter and Pointer]
\lean{Complexity.LogProg.access_inc_pointer}
\label{Complexity.LogProg.access_inc_pointer}
\leanok
\uses{Complexity.LogProg.AccessLabel, Complexity.LogProg.AccessPath, Complexity.LogProg.accessCfg, Complexity.LogProg.accessCfg_safe, Complexity.LogProg.access_inc, Complexity.LogProg.regCfg}
\difficulty{5}
Let $m$ be a Boolean mode, $y$ a Boolean input word, $r$ one of the four work tapes with
$r=0$ or $r=3$, $\ell$ an access label, $i,k,p$ natural numbers, and $q$ an integer.
Suppose $-1\le q\le$ the bit-length of $|y|+1$, and the value stored on tape $r$ in the
tuple $(i,0,k,p)$, plus one, is at most $|y|+1$. There is a finite sequence of
access-machine steps from the configuration with state “increment on tape $r$ with label
$\ell$,” tapes encoding $i,0,k,p$, and tape-$2$ head at $q$, to the configuration with
state “read label $\ell$ in phase eC,” with the counter for tape $0$ increased by one if
$r=0$, the counter for tape $3$ increased by one if $r=3$, and the other entries
unchanged; the tape-$2$ head remains at $q$.
\end{lemma}

\begin{lemma}[Returning from a restoring state]
\lean{Complexity.LogProg.access_return}
\label{Complexity.LogProg.access_return}
\leanok
\uses{Complexity.LogProg.AccessPath, Complexity.LogProg.AccessState, Complexity.LogProg.Dbl.o, Complexity.LogProg.accessCfg, Complexity.LogProg.accessCfg_safe, Complexity.LogProg.accessExit, Complexity.LogProg.accessPath_of_run, Complexity.LogProg.accessProg, Complexity.LogProg.accessTr, Complexity.LogProg.access_clear, Complexity.LogProg.backAct, Complexity.LogProg.length_bits_mono, Complexity.LogProg.regBack_x, Complexity.LogProg.regCfg, Complexity.LogProg.xCfg}
\difficulty{5}
Let $y$ be a Boolean word, let $eq$ be an optional Boolean value, and let $i,k,p,q$ be
natural numbers. If $i\le |y|+1$, $k\le |y|+1$, and $q\le |\mathrm{bits}(k)|$, then
there is a finite safe computation path from the restoring configuration with value
$eq$, parameters $i,k,p$, and tape-$2$ head position $q-1$ to the configuration
determined by the exit rule for $eq$, with parameters $0,k,p$.
\end{lemma}

\begin{definition}[Accessing and checking a list suffix]
\lean{Complexity.LogProg.accessTail}
\label{Complexity.LogProg.accessTail}
\leanok
Define a recursive operation on Boolean lists, with inputs a list $a$, a list $w$, a
Boolean $q$, and an optional Boolean $r$, whose value is an optional Boolean. If $a$ is
empty, the value is undefined when $r$ is false; otherwise it is the Boolean that holds
exactly when $q$ is true and $w$ is empty. If $a=b::a'$, the value is the result of
applying the operation to $a'$, the tail of $w$, the conjunction of $q$ with the
condition that the head of $w$ equals $b$, and the option containing $b$.
\end{definition}

\begin{lemma}[Tail comparison formula]
\lean{Complexity.LogProg.accessTail_eq}
\label{Complexity.LogProg.accessTail_eq}
\leanok
\uses{Complexity.LogProg.accessTail}
\difficulty{4}
Let $w_i$ and $w_r$ be finite bit strings, let $e$ be a Boolean, and let $a$ be either a
bit or absent. The tail comparison is absent if the last bit of $w_i$, when it exists,
and otherwise $a$, is $0$; in every other case it is present and equals $1$ precisely
when $e$ is true and $w_r=w_i$.
\end{lemma}

\begin{lemma}[Canonical access-tail result]
\lean{Complexity.LogProg.accessTail_canonical}
\label{Complexity.LogProg.accessTail_canonical}
\leanok
\uses{Complexity.LogProg.Canon, Complexity.LogProg.accessTail, Complexity.LogProg.accessTail_eq}
\difficulty{5}
Let $w_i$ and $w_r$ be finite bit strings, and let $e$ be a Boolean. The access-tail
operation on $w_i$ and $w_r$, initialized with Boolean state $e$ and no remembered last
bit, returns no value if $w_i$ is nonempty and ends in $0$; otherwise it returns the
Boolean conjunction of $e$ and the statement that $w_i=w_r$.
\end{lemma}

\begin{definition}[Access check]
\lean{Complexity.LogProg.accessCheck}
\label{Complexity.LogProg.accessCheck}
\leanok
\uses{Complexity.LogProg.Canon, Turing.pairDecode}
Given a Boolean list $y$ and a natural number $k$, decode $y$ as a pair of Boolean lists
if possible. If the second list in the decoded pair is canonical, return whether it
equals the binary representation of $k$; otherwise, or if decoding fails, return none.
\end{definition}

\begin{lemma}[Access check on encoded pairs]
\lean{Complexity.LogProg.accessCheck_pair}
\label{Complexity.LogProg.accessCheck_pair}
\leanok
\uses{Complexity.LogProg.accessCheck, Complexity.LogProg.bits_injective, Complexity.LogProg.canon_bits, Turing.pairDecode_pairEncode, Turing.pairEncode}
\difficulty{1}
For bit strings $z$ and natural numbers $r,k$, checking the encoding of $z$ paired with
the binary representation of $r$ returns a Boolean indicating whether $r=k$: the result
is `some true` exactly when $r=k$, and `some false` otherwise.
\end{lemma}

\begin{lemma}[Characterizing successful access checks]
\lean{Complexity.LogProg.accessCheck_some}
\label{Complexity.LogProg.accessCheck_some}
\leanok
\uses{Complexity.LogProg.accessCheck, Complexity.LogProg.bitsVal, Complexity.LogProg.bits_injective, Complexity.LogProg.canon_eq_bits, Turing.eq_pairEncode_of_pairDecode, Turing.pairDecode, Turing.pairEncode}
\difficulty{5}
If checking a bit string $y$ against a natural number $k$ returns `some b`, then there
are a bit string $z$ and a natural number $r$ such that $y$ is the encoding of $z$
paired with the binary representation of $r$, and $b$ is true exactly when $r=k$.
\end{lemma}

\begin{lemma}[A State Transition Step]
\lean{Complexity.LogProg.access_go}
\label{Complexity.LogProg.access_go}
\leanok
\uses{Complexity.LogProg.AccessState, Complexity.LogProg.accessCfg, Complexity.LogProg.goAct, Turing.Action.apply, Turing.moveInputPos_zero}
\difficulty{1}
Let $m$ be a Boolean mode, $y$ a Boolean input word, $s,s'$ access-machine states, and
$i,k,p$ natural numbers and $q$ an integer. Suppose that for every scanned input symbol
and every tuple of work-tape symbols, the transition prescribed in state $s$ is the
action that leaves the heads and tapes unchanged and moves to state $s'$. Then the
initial configuration with state $s$, tape contents encoding $i,0,k,p$, third work-tape
head at $q$, and all other work-tape heads at $0$ reaches the corresponding
configuration with state $s'$ in a finite sequence of access-machine steps.
\end{lemma}

\begin{lemma}[Work tape head movement]
\lean{Complexity.LogProg.access_move}
\label{Complexity.LogProg.access_move}
\leanok
\uses{Complexity.LogProg.AccessState, Complexity.LogProg.accessCfg, Complexity.LogProg.apply_xAct, Complexity.LogProg.xAct, Complexity.LogProg.xCfg, Turing.Action.apply, Turing.moveInputPos_zero}
\difficulty{3}
For any Boolean input word $y$, states $st,st'$, natural numbers $i,k,p$, integer head
position $q$, and sign $mv\in\{-1,0,1\}$, applying the action that leaves the input head
fixed, moves work tape $2$ by $mv$, and enters state $st'$ to the configuration in state
$st$ with head position $q$ on tape $2$ gives the configuration in state $st'$ with head
position $q+mv$ on tape $2$.
\end{lemma}

\begin{lemma}[Symbol read from the counter tape]
\lean{Complexity.LogProg.access_symbol}
\label{Complexity.LogProg.access_symbol}
\leanok
\uses{Complexity.LogProg.AccessState, Complexity.LogProg.accessCfg, Turing.Cfg.workTapeSymbols, Turing.FinTM.bufferTape}
\difficulty{1}
For any Boolean word $y$, access state $st$, and natural numbers $i,k,p,q$, the symbol
under the head of work tape $2$ in the corresponding access configuration is the first
bit, if any, of the binary expansion of $k$ after its first $q$ bits are dropped.
\end{lemma}

\begin{lemma}[Tail scan reaches its exit state]
\lean{Complexity.LogProg.access_tail_run}
\label{Complexity.LogProg.access_tail_run}
\leanok
\uses{Complexity.LogProg.AccessPath, Complexity.LogProg.accessCfg, Complexity.LogProg.accessCfg_safe, Complexity.LogProg.accessExit, Complexity.LogProg.accessReg, Complexity.LogProg.accessTail, Complexity.LogProg.accessTr, Complexity.LogProg.access_inc_pointer, Complexity.LogProg.access_move, Complexity.LogProg.access_read, Complexity.LogProg.access_return, Complexity.LogProg.access_step, Complexity.LogProg.access_symbol, Complexity.LogProg.length_bits_mono}
\difficulty{6}
Let $y$ and $w$ be finite bit strings, let $i,k,p,q$ be natural numbers, let $e$ be a
Boolean, and let $	ext{last}$ be either a bit or absent. Suppose $y$ with its first $i$
bits removed is $w$, $i\le |y|$, $k\le |y|+1$, and $q\le |\mathrm{bits}(k)|$, where
$\mathrm{bits}(k)$ is the binary word for $k$. Starting from the configuration with
input $y$, input head at position $1$, work tapes storing the binary words for
$i,0,k,p$, and head positions $0,0,q,0$, there is a finite safe computation in the
access machine (with parameter $\mathrm{mode}$) to the configuration whose state is the
exit determined by recursively scanning $w$ against the binary word for $k$ with its
first $q$ bits removed, updating the Boolean comparison and last-bit data according to
the scan; the final work tapes store the binary words for $0,0,k,p$ and all head
positions are zero.
\end{lemma}

\begin{lemma}[Rejecting from a restoring transition]
\lean{Complexity.LogProg.access_reject}
\label{Complexity.LogProg.access_reject}
\leanok
\uses{Complexity.LogProg.AccessPath, Complexity.LogProg.AccessState, Complexity.LogProg.accessCfg, Complexity.LogProg.accessCfg_safe, Complexity.LogProg.accessTr, Complexity.LogProg.access_move, Complexity.LogProg.access_return, Complexity.LogProg.access_step, Complexity.LogProg.xAct}
\difficulty{3}
Let $y$ be a Boolean word, $st$ an access state, and $i,k,p$ natural numbers with $i\le
|y|+1$ and $k\le |y|+1$. If, for every optional Boolean symbol $a$ and every work-head
reading function $w$, the transition from $st$ is the action that moves work tape $2$
one cell left and enters the restoring state with no equality result, then there is a
finite safe computation path from the configuration with state $st$ and parameters
$i,k,p$ to the configuration with rejecting answer and parameters $0,k,p$.
\end{lemma}

\begin{lemma}[A State Transition Step]
\lean{Complexity.LogProg.access_jump}
\label{Complexity.LogProg.access_jump}
\leanok
\uses{Complexity.LogProg.AccessPath, Complexity.LogProg.AccessState, Complexity.LogProg.accessCfg, Complexity.LogProg.accessCfg_safe, Complexity.LogProg.accessTr, Complexity.LogProg.access_go, Complexity.LogProg.access_step, Complexity.LogProg.goAct}
\difficulty{3}
Let $m$ be a Boolean mode, $y$ a Boolean input word, $s,s'$ access-machine states, and
$i,k,p$ natural numbers. Suppose that for every scanned input symbol and every tuple of
work-tape symbols, the access-machine transition in state $s$ is the stationary action
that moves to state $s'$. Then the configuration with state $s$, tapes encoding
$i,0,k,p$, and all work-tape heads at $0$ reaches the corresponding configuration with
state $s'$ in a finite sequence of access-machine steps.
\end{lemma}

\begin{lemma}[Invariance Under a Doubled Prefix]
\lean{Complexity.LogProg.accessCheck_dup}
\label{Complexity.LogProg.accessCheck_dup}
\leanok
\uses{Complexity.LogProg.accessCheck, Turing.pairDecode}
\difficulty{4}
For Boolean words $w$ and any natural number $k$, the access check on the word formed by
prefixing $w$ with two copies of the same bit $b$ equals the access check on $w$.
\end{lemma}

\begin{lemma}[Scanning a suffix of the input]
\lean{Complexity.LogProg.access_scan_run}
\label{Complexity.LogProg.access_scan_run}
\leanok
\uses{Complexity.LogProg.AccessPath, Complexity.LogProg.accessCfg, Complexity.LogProg.accessCheck, Complexity.LogProg.accessCheck_dup, Complexity.LogProg.accessExit, Complexity.LogProg.accessReg, Complexity.LogProg.accessTail_canonical, Complexity.LogProg.accessTr, Complexity.LogProg.access_inc_pointer, Complexity.LogProg.access_jump, Complexity.LogProg.access_read, Complexity.LogProg.access_reject, Complexity.LogProg.access_tail_run, Turing.pairDecode}
\difficulty{6}
Let $y$ and $w_i$ be Boolean words, and let $i,k,p$ be natural numbers. If the suffix of
$y$ obtained by dropping its first $i$ bits is $w_i$, $i\le |y|$, and $k\le |y|+1$, then
there is a finite safe computation path from the scan-reading configuration with
parameters $i,k,p$ to the configuration determined by the exit rule applied to the
access check of $w_i$ and $k$, with parameters $0,k,p$.
\end{lemma}

\begin{lemma}[Access counter increment path]
\lean{Complexity.LogProg.access_bump}
\label{Complexity.LogProg.access_bump}
\leanok
\uses{Complexity.LogProg.AccessPath, Complexity.LogProg.AccessState, Complexity.LogProg.Dbl.o, Complexity.LogProg.accessCfg, Complexity.LogProg.accessCfg_safe, Complexity.LogProg.accessPath_of_run, Complexity.LogProg.accessProg, Complexity.LogProg.access_frame, Complexity.LogProg.access_inc_pointer, Complexity.LogProg.inc_run, Complexity.LogProg.length_bits_mono, Complexity.LogProg.regCfg, Turing.FinTM.bufferTape}
\difficulty{6}
For either access mode, any bit-string input $y$ and natural numbers $k,p$ satisfying
$k+1\le |y|+1$ and $p+1\le |y|+1$ admit an access path from the configuration in the
first bump state, with four work tapes containing the binary representations of
$0,0,k,p$ and all work heads at position $0$, to the configuration in the initial scan
state with work tapes containing the binary representations of $0,0,k+1,p+1$ and all
work heads at position $0$.
\end{lemma}

\begin{definition}[Access value]
\lean{Complexity.LogProg.accessValue}
\label{Complexity.LogProg.accessValue}
\leanok
For a Boolean mode, a Boolean list $z$, and a natural number $r$, define the value to be
true precisely when $r<|z|$ if the mode is true; if the mode is false, use the $r$th
entry of $z$, defaulting to false when $r$ is out of range.
\end{definition}

\begin{definition}[Completed access configuration]
\lean{Complexity.LogProg.accessDone}
\label{Complexity.LogProg.accessDone}
\leanok
\uses{Complexity.LogProg.AccessState, Complexity.LogProg.accessCfg, Turing.Cfg}
For a Boolean answer $b$ and natural numbers $k$ and $p$, take the configuration
produced by the access procedure in its answer state with parameters $0,k,p$, then set
its machine state to the halting state and its output tape to the one-symbol list $[b]$.
\end{definition}

\begin{lemma}[Access path from an answer state]
\lean{Complexity.LogProg.access_answer}
\label{Complexity.LogProg.access_answer}
\leanok
\uses{Complexity.LogProg.AccessPath, Complexity.LogProg.accessCfg, Complexity.LogProg.accessCfg_safe, Complexity.LogProg.accessDone, Complexity.LogProg.accessTr, Complexity.LogProg.access_step, Complexity.LogProg.retAct, Turing.Action.apply, Turing.moveInputPos_zero}
\difficulty{3}
For either access mode, any bit-string input $y$, Boolean $b$, and natural numbers
$k,p$, there is an access path from the configuration in answer state $b$ whose four
work tapes contain the binary representations of $0,0,k,p$, respectively, and whose
third work head is at position $q=0$, to the same configuration with the machine halted
and output equal to the one-bit string $[b]$.
\end{lemma}

\begin{lemma}[Reading a Data Bit]
\lean{Complexity.LogProg.access_data_run}
\label{Complexity.LogProg.access_data_run}
\leanok
\uses{Complexity.LogProg.AccessPath, Complexity.LogProg.accessCfg, Complexity.LogProg.accessReg, Complexity.LogProg.accessTr, Complexity.LogProg.access_inc_pointer, Complexity.LogProg.access_jump, Complexity.LogProg.access_read, Turing.getElem?_dbl, Turing.length_dbl, Turing.length_pairEncode, Turing.pairEncode, Turing.pairEncode_eq_dbl}
\difficulty{5}
Let $m,e\in\{\mathrm{false},\mathrm{true}\}$, let $z,w$ be Boolean words, and let $k\le
|z|$. On input the self-delimiting pair encoding of $z$ and $w$, there is a finite
sequence of access-machine steps from the configuration in state “read data with
parameter $e$,” whose four tapes contain the binary words for $0,0,k,2k$ and whose third
work-tape head is at $2k$, to a configuration whose state and tapes are determined as
follows: if $k<|z|$ and $e$ is true, the state is “answer” with value $m\lor z_k$; if
$k<|z|$ and $e$ is false, the state is “bump false”; and if $k=|z|$, the state is
“answer false.” The final configuration has tape contents encoding $0,0,k,2k+1$, with
the third work-tape head at $2k+1$; the other work-tape heads are at $0$.
\end{lemma}

\begin{lemma}[Access computation returns the indexed bit]
\lean{Complexity.LogProg.access_valid}
\label{Complexity.LogProg.access_valid}
\leanok
\uses{Complexity.LogProg.AccessPath, Complexity.LogProg.accessCfg, Complexity.LogProg.accessCheck_pair, Complexity.LogProg.accessDone, Complexity.LogProg.accessValue, Complexity.LogProg.access_answer, Complexity.LogProg.access_bump, Complexity.LogProg.access_data_run, Complexity.LogProg.access_scan_run, Turing.length_pairEncode, Turing.pairEncode}
\difficulty{6}
Let $z$ be a finite bit string, $r$ a natural number, and $\mathrm{mode}$ a Boolean. For
every natural number $k$ with $k\le |z|$ and $k\le r$, there are natural numbers $k'$
and $p'$ such that, from the configuration on input $\langle z,\mathrm{bits}(r)\rangle$
with the access machine in its initial scan state, there is a finite safe computation to
a halting configuration that outputs the bit $\mathrm{mode}\land(r<|z|)$ if
$\mathrm{mode}$ is true, and otherwise outputs the $r$-th bit of $z$, using false when
that position is beyond the end of $z$.
\end{lemma}

\begin{definition}[Access language]
\lean{Complexity.LogProg.accessLang}
\label{Complexity.LogProg.accessLang}
\leanok
\uses{Complexity.LogProg.accessValue, Complexity.indexLang}
For a Boolean mode, define a language over Boolean strings by including exactly those
strings selected by the index-language construction whose access predicate, evaluated at
that mode and the corresponding indices, is true.
\end{definition}

\begin{lemma}[Acceptance computation for the access machine]
\lean{Complexity.LogProg.access_run}
\label{Complexity.LogProg.access_run}
\leanok
\uses{Complexity.LogProg.AccessPath, Complexity.LogProg.accessCfg, Complexity.LogProg.accessCheck, Complexity.LogProg.accessCheck_some, Complexity.LogProg.accessDone, Complexity.LogProg.accessLang, Complexity.LogProg.accessValue, Complexity.LogProg.access_answer, Complexity.LogProg.access_scan_run, Complexity.LogProg.access_valid, Complexity.LogProg.indicator_eq_decide, Complexity.indicator_indexLang_pairEncode, Turing.Cfg.init, Turing.MultiTapeTM.indicator, Turing.pairEncode}
\difficulty{5}
For every Boolean mode and Boolean word $y$, there exist natural numbers $k,p$ such that
the initial configuration on input $y$, in the scan-reading state, has a finite safe
computation path to a halting configuration whose output is the indicator of whether $y$
belongs to the language of pairs $(z,r)$ satisfying the access-value predicate for that
mode.
\end{lemma}

\begin{lemma}[Logarithmic-Space Access Language]
\lean{Complexity.LogProg.access_mem}
\label{Complexity.LogProg.access_mem}
\leanok
\uses{Complexity.LOGSPACE, Complexity.LogProg.Dbl.nilTM, Complexity.LogProg.Dbl.o, Complexity.LogProg.accessCfg, Complexity.LogProg.accessDone, Complexity.LogProg.accessLang, Complexity.LogProg.accessPath_run, Complexity.LogProg.accessProg, Complexity.LogProg.access_run, Complexity.LogProg.compileFinTM, Complexity.LogProg.compile_space, Complexity.LogProg.length_bits_le_log, Complexity.LogProg.log_poly_bound, Complexity.LogProg.rrun, Complexity.logSpace, Turing.Cfg.init, Turing.MultiTapeTM.indicator}
\difficulty{5}
For either Boolean mode $m$, the language consisting of pair encodings of Boolean words
$z$ and binary representations of natural numbers $r$ for which the access value of $z$
at $r$ in mode $m$ is true is decidable by a deterministic Turing machine using at most
a constant multiple of logarithmic space in the input length.
\end{lemma}

\begin{theorem}[Implicit logspace computability of identity]
\lean{Complexity.implicitlyLogspaceComputable_id}
\label{Complexity.implicitlyLogspaceComputable_id}
\leanok
\uses{Complexity.ImplicitlyLogspaceComputable, Complexity.LogProg.accessLang, Complexity.LogProg.accessValue, Complexity.LogProg.access_mem}
\difficulty{4}
The identity function on finite bit strings is implicitly logspace computable: its
output length is bounded by $C(|x|+1)^c$ for some natural numbers $C,c$ on every input
$x$, and both the language of encoded pairs $(x,	ext{bits}(i))$ for which output bit $i$
is $1$ and the language of such pairs for which $i$ is less than the output length are
decidable in logarithmic space.
\end{theorem}

\begin{theorem}[Input Bit Language Is Logarithmic-Space Decidable]
\lean{Complexity.inputBitLang_mem}
\label{Complexity.inputBitLang_mem}
\leanok
\uses{Complexity.LOGSPACE, Complexity.implicitlyLogspaceComputable_id, Complexity.inputBitLang}
\difficulty{1}
The language of binary strings of the form $\langle x,\operatorname{bits}(i)\rangle$,
where $\langle\cdot,\cdot\rangle$ is the encoding that doubles each bit of the first
string and then appends $01$ and the second string, and where the zero-indexed bit $x_i$
exists and equals $1$, is decidable by a finite-state-set Turing machine using at most
$c(\lfloor\log_2 n\rfloor+1)$ work-tape cells on inputs of length $n$, for some natural
number constant $c$.
\end{theorem}

\begin{theorem}[Input Length Language Is Logarithmic-Space Decidable]
\lean{Complexity.inputLenLang_mem}
\label{Complexity.inputLenLang_mem}
\leanok
\uses{Complexity.LOGSPACE, Complexity.implicitlyLogspaceComputable_id, Complexity.inputLenLang}
\difficulty{1}
The language of binary strings of the form $\langle x,\operatorname{bits}(i)\rangle$,
with the pairing encoding that doubles the bits of $x$ and then appends $01$ and the
bits of $i$, where $i<|x|$, is decidable by a finite-state-set Turing machine using at
most $c(\lfloor\log_2 n\rfloor+1)$ work-tape cells on inputs of length $n$, for some
natural number constant $c$.
\end{theorem}
